\documentclass[11pt,a4paper]{article}
\usepackage[margin=1in]{geometry}
\usepackage{times}
\usepackage{booktabs}
\usepackage{amsmath}
\usepackage{amssymb}
\usepackage{graphicx}
\usepackage{multirow}
\usepackage{xcolor}
\usepackage{hyperref}

\title{Cardinality-Constrained KV Cache Quantization:\\Preventing 2-Bit Compounding Failures\\on Hardware-Aligned Bit-Widths}

\author{Anonymous}

\begin{document}
\maketitle

%% ============================================================
%%  ABSTRACT
%% ============================================================
\begin{abstract}
Per-head mixed-precision KV cache quantization promises superior quality--compression trade-offs over uniform bit-width allocation.
However, on hardware-aligned non-consecutive bit-widths \{2,4,8,16\}, we discover a \emph{universal failure mode}: 2-bit quantization distortions compound non-additively through the residual stream, causing \emph{all} tested unconstrained solvers---greedy and MCKP---to over-allocate 2-bit heads and underperform naive uniform 4-bit across 8/8 diverse architectures (degradation 1.07$\times$--3.31$\times$ baseline).
A simple cardinality constraint ($\le k$ heads at 2-bit) universally prevents this over-allocation and, surprisingly, makes solver choice irrelevant: constrained greedy and constrained MCKP produce \emph{identical} bit allocations ($\Delta$=0.0000 PPL).
Under moderate compression (4$\times$), only 1/8 architectures benefits from non-uniform allocation; the remaining 7/8 should use uniform 4-bit.
A two-metric diagnostic (concentration $\gtrsim$0.50 \textsc{and} $t_1/t_5 \lesssim$0.70) correctly classifies all 9 tested models.
At 8K context length, uniform 4-bit degrades 11.15$\times$ baseline while constrained MCKP degrades only 1.23$\times$---a 9$\times$ gap in degradation rate---demonstrating critical deployment value.
\end{abstract}

%% ============================================================
%%  1. INTRODUCTION
%% ============================================================
\section{Introduction}
\label{sec:intro}

KV cache memory is the primary bottleneck for long-context LLM inference.
Per-head mixed-precision quantization allocates different bit-widths to different attention heads based on their sensitivity, promising better quality--compression trade-offs than uniform allocation.
Prior work such as RateQuant~\cite{ratequant2024} provides rate-distortion optimal allocation on consecutive integer bit-widths \{2,\ldots,6\}, leveraging polymatroid structure for provably optimal greedy solutions.

However, practical deployment requires \emph{hardware-aligned} bit-widths \{2,4,8,16\} supported by efficient GPU kernels.
On these non-consecutive bit-widths, we discover that:
\begin{enumerate}
\item \textbf{2-bit compounding is universal.} Independent per-head profiling underestimates joint distortion when multiple heads are simultaneously quantized to 2-bit. Per-head profiling predicts joint PPL$\approx$11 but actual joint 2-bit PPL is 32--37 (3$\times$ prediction error). All unconstrained solvers consequently over-allocate 2-bit heads and underperform uniform 4-bit across 8/8 tested architectures.

\item \textbf{A cardinality constraint is the universal fix---and makes solver choice irrelevant.} Capping the number of 2-bit heads at $k\le 8$ prevents compounding across all architectures. Under this constraint, greedy and MCKP produce identical allocations ($\Delta$=0.0000 PPL), eliminating solver complexity.

\item \textbf{Non-uniform benefit is rare.} Under moderate (4$\times$) compression, only 1/8 architectures---those with a ``spread'' sensitivity pattern---benefits from non-uniform allocation. A two-metric diagnostic correctly predicts this for all 9 tested models, providing practitioners a simple check before investing in profiling.
\end{enumerate}

Our contributions are: (i) the first systematic characterization of 2-bit compounding effects in per-head KV cache quantization, (ii) a universal constraint fix that eliminates solver complexity, and (iii) the first empirical study of \emph{when} non-uniform allocation is beneficial, correcting the implicit assumption that it always helps.

%% ============================================================
%%  2. RELATED WORK
%% ============================================================
\section{Related Work}
\label{sec:related}

\paragraph{KV Cache Quantization.}
KIVI~\cite{kivi2024} applies per-channel group quantization with residual storage, achieving near-lossless 4-bit compression.
Our work is orthogonal: KIVI optimizes \emph{within-bit} quantization quality via group granularity, while we optimize \emph{cross-bit} per-head allocation and characterize its failure modes.
The two approaches are complementary and can be combined.

\paragraph{Rate-Distortion Allocation.}
RateQuant~\cite{ratequant2024} formulates per-head allocation as rate-distortion optimization, proving greedy optimality on consecutive integers \{2,\ldots,6\} via polymatroid structure (Proposition~6).
On hardware-aligned \{2,4,8,16\}, this structure breaks: marginal gain per bit between 2$\to$4 and 4$\to$8 is non-monotone, and greedy over-allocates 2-bit heads (68/112 = 60.7\% on Qwen, PPL 32.45 vs optimal $\sim$10.91).
RateQuant acknowledges compounding effects (Appendix~M) but does not address them.
Our cardinality constraint fixes RateQuant's greedy, achieving identical results to MCKP.

\paragraph{Mixed-Precision Quantization.}
Weight quantization literature extensively studies per-layer or per-group mixed-precision allocation.
KV cache quantization differs: attention heads exhibit architecture-specific sensitivity patterns that determine whether non-uniform allocation helps at all---a question not previously studied.

%% ============================================================
%%  3. METHOD
%% ============================================================
\section{Method}
\label{sec:method}

\subsection{Problem Formulation}

Given a model with $H$ KV heads and bit-width options $\mathcal{B}=\{2,4,8,16\}$, we seek an assignment $b_h \in \mathcal{B}$ for each head $h$ minimizing total distortion subject to a memory budget:
\begin{equation}
\min_{\{b_h\}} \sum_{h=1}^{H} d_h(b_h) \quad \text{s.t.} \quad \sum_{h=1}^{H} b_h \le B
\end{equation}
where $d_h(b_h)$ is the per-head distortion measured by the change in perplexity when head $h$ is quantized to $b_h$ bits (all other heads at full precision).

This is a standard Multiple-Choice Knapsack Problem (MCKP).
Independent per-head profiling assumes distortions are additive: $D(\{b_h\}) = \sum_h d_h(b_h)$.
We show this assumption fails when multiple heads are assigned 2-bit simultaneously.

\subsection{The 2-Bit Compounding Effect}

When $k$ heads are simultaneously quantized to 2-bit, the joint distortion exceeds the sum of individual distortions:
\begin{equation}
D_{\text{joint}}(k) \gg \sum_{h \in S_k} d_h(2)
\label{eq:nonadditivity}
\end{equation}
We characterize this non-additivity empirically via a compounding sweep: varying $k$ from 0 to $H$ while measuring actual PPL.
On Qwen2.5-7B, profiling predicts joint PPL$\approx$11 at $k$=112 but actual PPL is 98.75 (9$\times$ prediction error).

\subsection{Cardinality-Constrained Allocation}

We augment the MCKP with a cardinality constraint:
\begin{equation}
|\{h : b_h = 2\}| \le k
\label{eq:constraint}
\end{equation}
This constraint limits the number of 2-bit heads, preventing compounding.
The constrained problem remains MCKP-solvable.
We determine $k$ via a small-scale compounding test: sweep $k$ on a calibration set and select the largest $k$ with $<2.5\%$ degradation.
Empirically, $k \le 8$ is universally safe across all tested architectures.

\subsection{Sensitivity Distribution Diagnostic}

To determine whether non-uniform allocation is beneficial \emph{a priori}, we propose a two-metric diagnostic based on the per-head sensitivity profile $\{d_h(2)\}$:

\begin{itemize}
\item \textbf{Top-5\% concentration}: fraction of total sensitivity in the top 5\% of heads. Threshold: $\gtrsim 0.50$.
\item \textbf{Dominance ratio} $t_1/t_5$: ratio of the single most sensitive head to the sum of the top-5 heads. Threshold: $\lesssim 0.70$.
\end{itemize}

The joint condition (concentration $\gtrsim 0.50$ \textsc{and} $t_1/t_5 \lesssim 0.70$) identifies ``spread'' patterns where multiple heads have high but differentiated sensitivity, enabling MCKP to exploit cross-bitwidth gradients.
A ``spike'' pattern (high concentration but high $t_1/t_5$) indicates a single dominant outlier that cannot be exploited.
A ``flat'' pattern (low concentration) indicates homogeneous sensitivity where uniform allocation is optimal.

%% ============================================================
%%  4. EXPERIMENTS
%% ============================================================
\section{Experiments}
\label{sec:experiments}

\subsection{Setup}

We evaluate on 8 architectures spanning GQA ratios from 1:1 to 8:1: Qwen2.5-7B (7:1, 112 KV heads), Yi-1.5-9B (8:1, 192 heads), Llama-3.1-8B (4:1, 256 heads), Mistral-7B-v0.3 (4:1, 256 heads), Gemma-2-9B (2:1, 336 heads), Pythia-1.4B (MHA, 384 heads), Qwen2.5-14B (7:1, 384 heads), and Phi-3-medium (400 heads).
Profiling and evaluation use WikiText-2 at 4$\times$ compression (budget=0.25) unless stated otherwise.
High-resolution profiling uses 64 sequences $\times$ 2048 tokens.

\subsection{Unconstrained Solvers Universally Fail}
\label{sec:unconstrained}

Table~\ref{tab:main} shows that unconstrained allocation (both greedy and MCKP) underperforms uniform 4-bit across all tested architectures.
Degradation follows a concentration-dependent gradient:
\emph{catastrophic} on spread models (Qwen 3.31$\times$),
\emph{moderate} on flat-high-2bit models (Gemma-2 1.47$\times$, where the solver assigns 106/336=31.5\% of heads to 2-bit because distortions are indistinguishable),
and \emph{mild} on spike/flat-low-2bit models (Yi 1.08$\times$, Llama 1.07$\times$, Mistral 1.08$\times$).

\begin{table}[t]
\centering
\small
\caption{Main results across architectures at 4$\times$ compression. \textbf{Unconstr.}: MCKP without cardinality constraint. \textbf{Constr.}: MCKP with $k$=8 (or natural $k$ where noted). \textbf{Greedy}: RateQuant Algorithm~1 with constraint. All unconstrained solvers underperform uniform 4-bit.}
\label{tab:main}
\begin{tabular}{@{}lcccccc@{}}
\toprule
\textbf{Model} & \textbf{BF16} & \textbf{Uniform} & \textbf{Constr.} & \textbf{Greedy} & \textbf{Unconstr.} & \textbf{Shape} \\
 & & \textbf{4-bit} & \textbf{MCKP} & \textbf{Constr.} & \textbf{MCKP} & \\
\midrule
Qwen-7B & 9.81 & 15.45 & \textbf{10.91} & \textbf{10.91}\textsuperscript{$\dagger$} & 36.92 & Spread \\
Llama-8B & 8.09 & 8.47 & 8.45 & --- & 8.63 & Spike \\
Mistral-7B & 7.04 & 7.13 & 7.16 & --- & 7.57 & Flat \\
Gemma-2-9B & 6.03 & 6.17 & 6.20 & --- & 8.85 & Flat \\
Yi-1.5-9B & 7.02 & 7.22 & 7.22 & --- & 7.58\textsuperscript{*} & Spike \\
Pythia-1.4B & 16.49 & 24.29 & 24.45 & --- & 31.54 & Flat \\
Qwen-14B & 6.87 & 7.20 & 7.22 & 7.22\textsuperscript{$\dagger$} & 7.69 & Flat \\
\bottomrule
\end{tabular}

\vspace{2pt}
{\footnotesize \textsuperscript{$\dagger$}Greedy constrained $=$ MCKP constrained ($\Delta$=0.0000). \textsuperscript{*}Estimated from 1.08$\times$ degradation ratio.}
\end{table}


\subsection{Solver Equivalence under Constraint}
\label{sec:solver_equiv}

Table~\ref{tab:solver_equiv} demonstrates that the cardinality constraint is the primary intervention, not the solver.
On Qwen-7B, unconstrained greedy (PPL 32.45) and unconstrained MCKP (PPL 36.92) both fail catastrophically.
With $k$=8, both produce \emph{identical} bit allocations (\{2:8, 4:100, 8:4\}) and identical PPL (11.2312).
This equivalence replicates on Qwen-14B ($\Delta$=0.0000).

\begin{table}[t]
\centering
\small
\caption{Solver equivalence on Qwen-7B (112 heads, 4$\times$ compression). The constraint eliminates solver choice entirely.}
\label{tab:solver_equiv}
\begin{tabular}{@{}lccc@{}}
\toprule
\textbf{Solver} & \textbf{PPL} & \textbf{2-bit heads} & \textbf{Distribution} \\
\midrule
Greedy unconstrained & 32.45 & 68 & \{2:68, 4:30, 8:4, 16:10\} \\
MCKP unconstrained & 36.92 & $\sim$70 & --- \\
\midrule
Greedy + $k$=8 & \textbf{11.2312} & 8 & \{2:8, 4:100, 8:4\} \\
MCKP + $k$=8 & \textbf{11.2312} & 8 & \{2:8, 4:100, 8:4\} \\
\midrule
Uniform 4-bit & 15.45 & 0 & all 4-bit \\
\bottomrule
\end{tabular}
\end{table}


\subsection{Compounding Characterization}
\label{sec:compounding}

\paragraph{Qwen Compounding Sweep.}
Table~\ref{tab:compounding} shows PPL as a function of the maximum number of 2-bit heads $k$ on Qwen2.5-7B.
The curve is U-shaped: $k$=8 is the sweet spot (2.4\% degradation), while $k$=64 causes 2.34$\times$ and $k$=112 causes 10.07$\times$ degradation.

\begin{table}[t]
\centering
\small
\caption{Compounding sweep on Qwen2.5-7B. PPL grows super-linearly beyond $k$=8.}
\label{tab:compounding}
\begin{tabular}{@{}rcc@{}}
\toprule
$k$ & PPL & Ratio to baseline \\
\midrule
0 & 9.81 & 1.00$\times$ \\
8 & 10.05 & 1.02$\times$ \\
16 & 10.42 & 1.06$\times$ \\
32 & 11.60 & 1.18$\times$ \\
64 & 22.95 & 2.34$\times$ \\
112 & 98.75 & 10.07$\times$ \\
\bottomrule
\end{tabular}
\end{table}

\paragraph{Cross-Architecture Compounding.}
Table~\ref{tab:cross_compound} shows that $k$=8 is universally safe ($<$2.5\% degradation), but curve shapes differ:
Qwen exhibits a U-shape (non-uniform benefit exists at the sweet spot),
Llama shows gentle monotonic increase,
and Mistral is flat-monotonic (every 2-bit head hurts).

\begin{table}[t]
\centering
\small
\caption{Cross-architecture compounding at $k$=8. All architectures show $<$2.5\% degradation, confirming $k$=8 as a universally safe threshold.}
\label{tab:cross_compound}
\begin{tabular}{@{}lcccc@{}}
\toprule
\textbf{Model} & \textbf{$k$=8 PPL} & \textbf{Baseline} & \textbf{Degrad.} & \textbf{Curve} \\
\midrule
Qwen-7B & 10.05 & 9.81 & 2.4\% & U-shaped \\
Llama-8B & 8.23 & 8.09 & 1.7\% & Monotonic \\
Mistral-7B & 7.06 & 7.04 & 0.4\% & Monotonic \\
\bottomrule
\end{tabular}
\end{table}

\paragraph{Top-$k$ vs.\ Random-$k$.}
To confirm non-additivity, we compare assigning the $k$ most sensitive heads to 2-bit (top-$k$) versus $k$ random heads (random-$k$) on Qwen (Table~\ref{tab:topk_vs_random}).
Both exhibit super-linear compounding, but random-$k$ is worse due to potential selection of correlated heads within the same layer.

\begin{table}[t]
\centering
\small
\caption{Top-$k$ vs.\ random-$k$ compounding on Qwen2.5-7B. Both exhibit super-linear growth; profiling-based selection (top-$k$) mitigates but does not eliminate compounding.}
\label{tab:topk_vs_random}
\begin{tabular}{@{}rcccc@{}}
\toprule
$k$ & Top-$k$ & Random-$k$ (mean$\pm$std) & Ratio (top) & Ratio (rand) \\
\midrule
8 & 10.04 & 11.07$\pm$0.69 & 1.56$\times$ & 1.76$\times$ \\
16 & 10.42 & 12.91$\pm$0.40 & 1.56$\times$ & 2.10$\times$ \\
32 & 11.60 & 18.33$\pm$2.53 & 1.78$\times$ & 2.73$\times$ \\
64 & 22.95 & 55.94$\pm$9.73 & 4.89$\times$ & 8.88$\times$ \\
\bottomrule
\end{tabular}
\end{table}


\subsection{When Does Non-Uniform Allocation Help?}
\label{sec:when}

Table~\ref{tab:diagnostic} shows the sensitivity distribution diagnostic across 9 models.
Only Qwen2.5-7B meets both conditions (concentration$\ge$0.50, $t_1/t_5 \le$0.70), correctly predicting it as the sole beneficiary of non-uniform allocation.
Yi-1.5 (concentration 0.71 but $t_1/t_5$=0.90) and CodeLlama (0.80, 0.81) illustrate why concentration alone is insufficient: single extreme outliers (spike pattern) do not provide optimization opportunities for MCKP.

\begin{table}[t]
\centering
\small
\caption{Sensitivity distribution diagnostic. The joint condition (conc.$\ge$0.50 \textsc{and} $t_1/t_5 \le$0.70) achieves 9/9 accuracy in predicting non-uniform benefit.}
\label{tab:diagnostic}
\begin{tabular}{@{}lccccc@{}}
\toprule
\textbf{Model} & \textbf{Conc.} & $t_1/t_5$ & \textbf{Shape} & \textbf{Predicted} & \textbf{Actual} \\
\midrule
Qwen-7B & 0.61 & 0.57 & Spread & Benefit & +27.3\% \\
Llama-8B & 0.57 & 0.85 & Spike & No benefit & +0.2\% \\
Yi-1.5-9B & 0.71 & 0.90 & Spike & No benefit & 0\% \\
Mistral-7B & 0.22 & --- & Flat & No benefit & $-$0.4\% \\
Gemma-2-9B & 0.23 & 0.38 & Flat & No benefit & $-$0.3\% \\
Pythia-1.4B & 0.24 & --- & Flat & No benefit & $-$0.7\% \\
Qwen-14B & 0.33 & --- & Flat & No benefit & $-$0.3\% \\
Phi-3-med & 0.17 & 0.24 & Diffuse & No benefit & ---\textsuperscript{$\star$} \\
CodeLlama & 0.80 & 0.81 & Spike & No benefit & ---\textsuperscript{$\star$} \\
\bottomrule
\end{tabular}

\vspace{2pt}
{\footnotesize \textsuperscript{$\star$}Profiling-only (full pipeline not run due to resource constraints); diagnostic prediction consistent with profiling pattern.}
\end{table}


\subsection{Budget Sweep}
\label{sec:budget}

Table~\ref{tab:budget} shows that compression ratio amplifies the benefit of constrained MCKP.
At 2$\times$, the budget is sufficient and differences are small.
At 4$\times$ (the primary operating point), MCKP achieves 29.4\% improvement over uniform.
At 8$\times$, uniform 2-bit is catastrophic (PPL 2644) while constrained MCKP achieves PPL 57.6---still high but 45$\times$ better.

\begin{table}[t]
\centering
\small
\caption{Budget sweep on Qwen2.5-7B. Compression amplifies non-uniform benefit.}
\label{tab:budget}
\begin{tabular}{@{}cccc@{}}
\toprule
\textbf{Compression} & \textbf{MCKP PPL} & \textbf{Uniform PPL} & \textbf{Improvement} \\
\midrule
2$\times$ & 9.85 & 10.20 & +3.4\% \\
4$\times$ & 10.91 & 15.45 & +29.4\% \\
8$\times$ & 57.60 & 2644 & $>$97\% \\
\bottomrule
\end{tabular}
\end{table}


\subsection{Long-Context Evaluation}
\label{sec:longctx}

Table~\ref{tab:longctx} demonstrates that long context amplifies the compounding failure.
At 8K tokens, uniform 4-bit degrades to 11.15$\times$ baseline (PPL 77.75) while constrained MCKP maintains 1.23$\times$ (PPL 8.59)---a 9$\times$ gap in degradation rate.

\begin{table}[t]
\centering
\small
\caption{Long-context evaluation on Qwen2.5-7B at 4$\times$ compression.}
\label{tab:longctx}
\begin{tabular}{@{}lcccccc@{}}
\toprule
\textbf{Ctx} & \textbf{BF16} & \textbf{MCKP} & \textbf{Uniform} & \textbf{MCKP ratio} & \textbf{Unif.\ ratio} \\
\midrule
2K & 9.81 & 10.91 & 15.45 & 1.11$\times$ & 1.57$\times$ \\
8K & 6.97 & 8.59 & 77.75 & 1.23$\times$ & 11.15$\times$ \\
\bottomrule
\end{tabular}
\end{table}


\subsection{Downstream Task Evaluation}
\label{sec:downstream}

Table~\ref{tab:downstream_qwen} shows that Qwen's PPL improvement translates to consistent downstream gains (+1.25 to +4.00 percentage points, 72.8\% average recovery).
Table~\ref{tab:downstream_cross} shows that the diagnostic correctly predicts downstream outcomes across all 5 evaluated models.
Non-spread architectures show minor systematic degradation (Yi $-$0.77pp, Llama $-$0.99pp, Mistral $-$0.48pp) rather than preservation.
ARC-Challenge shows the largest per-task sensitivity to allocation strategy, consistent with its reliance on multi-step reasoning where subtle KV cache errors accumulate (Yi $-$3.24pp, Llama $-$2.65pp).

\begin{table}[t]
\centering
\small
\caption{Downstream evaluation on Qwen2.5-7B (0-shot, 4$\times$ compression). MCKP recovers 72.8\% of the BF16-to-uniform gap on average.}
\label{tab:downstream_qwen}
\begin{tabular}{@{}lcccc@{}}
\toprule
\textbf{Task} & \textbf{BF16} & \textbf{Uniform} & \textbf{MCKP} & $\Delta$ (pp) \\
\midrule
ARC-C & 0.529 & 0.481 & 0.514 & +3.16 \\
HellaSwag & 0.783 & 0.728 & 0.768 & +4.00 \\
PIQA & 0.793 & 0.766 & 0.779 & +1.25 \\
WinoGrande & 0.690 & 0.662 & 0.682 & +1.97 \\
\bottomrule
\end{tabular}
\end{table}

\begin{table}[t]
\centering
\small
\caption{Cross-architecture downstream summary (0-shot average, 5 tasks). Diagnostic correctly predicts both PPL and downstream outcomes for all 5 models.}
\label{tab:downstream_cross}
\begin{tabular}{@{}lcccl@{}}
\toprule
\textbf{Model} & \textbf{Shape} & $\Delta$ \textbf{(pp)} & \textbf{Predicted} & \textbf{Correct} \\
\midrule
Qwen-7B & Spread & +1.92 & Benefit & \checkmark \\
Yi-1.5-9B & Spike & $-$0.77 & No benefit & \checkmark \\
Llama-8B & Spike & $-$0.99 & No benefit & \checkmark \\
Mistral-7B & Flat & $-$0.48 & No benefit & \checkmark \\
Gemma-2-9B & Flat & +0.14 & No benefit & \checkmark \\
\bottomrule
\end{tabular}
\end{table}

\paragraph{Multi-seed confidence intervals.}
Three calibration seeds on Qwen2.5-7B yield PPL 10.99$\pm$0.065 (CV=0.59\%), with identical macro allocation (2:8, 4:100, 8:4) across all seeds.
Downstream confidence intervals confirm robustness: HellaSwag +4.13pp (std=0.002), PIQA +1.96pp (std=0.002), WinoGrande +1.58pp (std=0.012).
ARC-Challenge shows +0.00pp (std=0.004), consistent with this task's higher sensitivity to allocation details.

\subsection{Ablation Studies}
\label{sec:ablations}

\paragraph{L0 Exclusion Removal.}
Removing the hard L0 exclusion rule from the solver produces identical PPL (10.91 in both cases): the profiling signal is strong enough for the solver to automatically assign L0's 4 sensitive heads to 8-bit without domain knowledge.

\paragraph{Profiling Resolution.}
High-resolution profiling (64$\times$2048) yields PPL 10.91 vs.\ low-resolution (16$\times$512) PPL 11.23, a 2.9\% improvement.
High-resolution distinguishes 100 heads vs.\ 37, but the macro bit distribution is identical (2:8, 4:100, 8:4), confirming the constraint's robustness to profiling quality.

\paragraph{K/V Separation.}
Treating K and V heads as independent MCKP items yields PPL 11.22 vs.\ joint K+V PPL 10.91.
The joint formulation is superior because budget sharing between K and V naturally discovers that V heads are more robust to quantization.

\paragraph{Oracle Comparison.}
Brute-force optimal allocation on 8 layer groups shows MCKP is exactly optimal on 5/8 groups, with maximum gap 0.018 PPL (0.18\% of baseline) on the remaining 3---position swaps within the same bit budget, confirming constrained MCKP is near-optimal.

\paragraph{Cross-Dataset Calibration.}
Calibrating on C4 instead of WikiText-2 yields 96.4\% allocation agreement (108/112 heads identical) and PPL 11.29 (vs.\ 10.91 with WikiText-2 calibration), confirming robustness to calibration domain.
Cross-domain transfer (WikiText-2$\to$C4 eval) maintains the MCKP advantage: 20.56 vs.\ 35.42 uniform.

\paragraph{Heuristic Baselines.}
Table~\ref{tab:heuristics} compares constrained MCKP against simpler allocation strategies.
MCKP consistently outperforms L0-constrained random allocation and naive positional assignment across all 3 models, confirming profiling-based selection has genuine value on architectures where non-uniform allocation is beneficial.

\begin{table}[t]
\centering
\small
\caption{Allocation baseline comparison across 3 architectures. MCKP outperforms all heuristics; heuristics still beat uniform only on Qwen (spread).}
\label{tab:heuristics}
\begin{tabular}{@{}lcccc@{}}
\toprule
\textbf{Model} & \textbf{MCKP} & \textbf{L0-rand (med)} & \textbf{Naive pos.} & \textbf{Uniform} \\
\midrule
Qwen-7B & \textbf{10.91} & 11.52 & 12.23 & 15.45 \\
Llama-8B & \textbf{8.45} & 10.04 & 9.94 & 8.47 \\
Mistral-7B & \textbf{7.16} & 7.59 & 7.48 & 7.13 \\
\bottomrule
\end{tabular}
\end{table}


\paragraph{Mistral $k$ Ablation.}
On Mistral, PPL increases monotonically with $k$: $k$=4$\to$7.165, $k$=8$\to$7.200, $k$=12$\to$7.245, $k$=16$\to$7.311, $k$=18$\to$7.323 (uniform=7.131).
Every additional 2-bit head hurts, confirming flat-distribution architectures should use uniform 4-bit.

\subsection{Runtime Analysis}
\label{sec:runtime}

Table~\ref{tab:runtime} shows the latency overhead of mixed-precision KV cache.
Constrained MCKP incurs 8.5\% decode latency overhead vs.\ uniform 4-bit due to per-head dynamic dequantization dispatch, with identical memory footprint.
The MCKP solver itself requires only 5ms (one-time cost).
At long context where the quality gap is largest (8K: 1.23$\times$ vs.\ 11.15$\times$ degradation), the latency cost is justified.

\begin{table}[t]
\centering
\small
\caption{Runtime benchmark on Qwen2.5-7B (seq\_len=2048).}
\label{tab:runtime}
\begin{tabular}{@{}lcccc@{}}
\toprule
\textbf{Method} & \textbf{Prefill} & \textbf{Decode} & \textbf{Throughput} & \textbf{Memory} \\
 & (ms) & (ms/tok) & (tok/s) & (MB) \\
\midrule
BF16 & 45.86 & 33.78 & 29.6 & 21 \\
Uniform 4-bit & 54.25 & 42.46 & 23.6 & 5.25 \\
Constr.\ MCKP & 57.03 & 46.09 & 21.8 & 5.25 \\
\bottomrule
\end{tabular}
\end{table}


\subsection{KIVI Comparison}
\label{sec:kivi}

KIVI~\cite{kivi2024} 4-bit with group size 32 and residual length 128 achieves PPL 9.90 on Qwen---better than our constrained MCKP (10.91).
However, the two approaches are orthogonal: KIVI optimizes per-channel group quantization granularity (within-bit quality), while we optimize per-head bit-width allocation (cross-bit allocation) and characterize its failure modes.
Combining KIVI's quantization with our allocation framework is a natural extension.

\begin{table}[t]
\centering
\small
\caption{KIVI comparison on Qwen2.5-7B (4$\times$ compression). Approaches are complementary.}
\label{tab:kivi}
\begin{tabular}{@{}lcc@{}}
\toprule
\textbf{Method} & \textbf{PPL} & \textbf{Approach} \\
\midrule
BF16 baseline & 9.81 & --- \\
KIVI 4-bit (g32, r128) & 9.90 & Per-channel \\
Constrained MCKP & 10.91 & Per-head \\
Uniform 4-bit (ours) & 15.45 & Per-head \\
\bottomrule
\end{tabular}
\end{table}


%% ============================================================
%%  5. CONCLUSION
%% ============================================================
\section{Conclusion}
\label{sec:conclusion}

We show that per-head KV cache quantization on hardware-aligned bit-widths \{2,4,8,16\} has a universal failure mode: 2-bit distortions compound non-additively, causing all unconstrained solvers to underperform uniform 4-bit across 8 architectures.
A cardinality constraint ($\le k$ heads at 2-bit) universally prevents this, making solver choice irrelevant (greedy $=$ MCKP, $\Delta$=0.0000).
Under moderate compression, only 1/8 architectures benefits from non-uniform allocation---most models should use uniform 4-bit, checkable via a simple two-metric diagnostic (9/9 accuracy).
At long context (8K), the constraint reduces degradation from 11.15$\times$ to 1.23$\times$ baseline.

\paragraph{Limitations.}
Our study covers 8 architectures up to 14B parameters; 72B-scale profiling was resource-limited. The diagnostic is validated on $n$=9 models---directional evidence, not a statistical claim. Runtime overhead (8.5\%) may vary with longer sequences. Multi-seed downstream confidence intervals are available only for Qwen.

\paragraph{Future Work.}
Combining cardinality-constrained allocation with KIVI-style group quantization, extending to token-level eviction policies, and validating on 70B+ models.

\begin{thebibliography}{9}
\bibitem{ratequant2024}
Zhihang Yuan, Yuzhang Shang, Yang Zhou, Zhen Dong, Chenhao Xue, Bingzhe Wu, Zhikai Li, Qingyi Gu, Yong Jae Lee, Yan Yan, Beidi Chen, Guangyu Sun, and Kurt Keutzer.
{RateQuant: Rate-Distortion Optimal KV Cache Quantization}.
\emph{arXiv preprint}, 2024.

\bibitem{kivi2024}
Zirui Liu, Jiayi Yuan, Hongye Jin, Shaochen Zhong, Zhaozhuo Xu, Vladimir Braverman, Beidi Chen, and Xia Hu.
{KIVI: A Tuning-Free Asymmetric 2bit Quantization for KV Cache}.
\emph{Proceedings of ICML}, 2024.
\end{thebibliography}

\end{document}
